\section{Hardness of batches}\label{sec:batches}

The two parts of Theorem~\ref{thm:gap} give two ways to recover the
count specified by one description. Wiring the signed processors realizes the
success count of a public program. Substituting a hypothetical batch
network into that wiring lets an evaluator recover the count with
enough time left, after compression, for the program to add one.

\begin{proof}[Proof of Theorem~\ref{thm:batches}]
Fix the finite field $F$ and write $p=\operatorname{char}F$.
Choose the compressor scalars in $\mathbb F_p$. Take $b,\rho$ as in
Theorem~\ref{thm:gap}, decreasing $\rho$ below as needed, and take
$0<\delta<b/2$ and $0<\sigma<\rho$. Suppose that for some sufficiently
large $n$ and some $1\le k\le 2^{n^\sigma}$ a trivalent network
$\mathcal B$ of volume at most $kn^{1+\delta}$ realizes
$(U_n^*)^{\otimes k}$ over $F$. Put
\[
                    R=k\,2^{\lceil n^\rho\rceil}.
\]
This choice meets the clock hypotheses and partitions the processors
into $R/k$ whole batches. Use the wiring of Theorem~\ref{thm:gap}(1)
with the private machine fixed by Theorem~\ref{thm:gap}(2).
For a description $D$, let $Z(D)\in\mathbb F_p$ be the value of this
wiring with its pins set to $D$. It is the integer success count of the
stored program, reduced modulo $p$.

Reserve fields in $D$ for its public code, the size parameters, and
an indexed description of $\mathcal B$ with all its local entries in
$F$. Each tensor has bounded degree and $F$ is fixed, so the network
uses $O(kn^{1+\delta})$ words, even giving each wire endpoint and each
local entry its own word. All active fields, including the code chosen
below, fit in $\lfloor R/10\rfloor$ words for large $n$. The public
machine can read its entire description, including its program.

\smallskip
\noindent\emph{Evaluate the substituted network.}
Fix an evaluator $E$ whose input is a description $D$ with this network
and these size parameters. Generate the pins and links for $D$, and
replace each consecutive group of $k$ processors by a copy of
$\mathcal B$. This preserves the scalar in $F$ for every choice of the
public code in $D$. Guess each summed wire bit once and multiply the
selected tensor entries in $F$, enforcing the pins and links.

To turn a product $u\in F$ into an acceptance count, fix an
$\mathbb F_p$-linear map $\lambda:F\to\mathbb F_p$ with $\lambda(1)=1$.
Guess $\lceil\log p\rceil$ bits and accept when the resulting integer
is smaller than the representative of $\lambda(u)$ in $[0,p-1]$.
Summing over assignments applies $\lambda$ to the entire contraction.
Since $Z(D)$ lies in $\mathbb F_p$, the evaluator satisfies
\[
                    \#E(D)\equiv Z(D)\pmod p ,
\]
where $\#E(D)$ is its integer number of accepting paths and the
congruence identifies $\mathbb F_p$ with residues modulo $p$.

The evaluator only reads the stored instructions as literal data
when setting the pins. Generating the genuine wiring takes
$O(Rn\operatorname{poly}(w))$ ticks by Theorem~\ref{thm:gap}(1).
Each exterior slot has an $O(w)$-bit identifier for its batch and
corresponding leg. Traversing the $R/k$ copies of $\mathcal B$,
checking the assignments and multiplying their entries therefore takes
$O(Rn^{1+\delta}\operatorname{poly}(w))$ ticks in all. The exponent
on $w$ is fixed, so $\delta<b/2$ and a sufficiently small $\rho$ put
this procedure within $T_+$, without private jobs.

\smallskip
\noindent\emph{Choose the program whose pins are evaluated.}
Give the public program two branches. The first accepts once.
The second runs the compressed simulation of $E$ on its actual $D$
using Theorem~\ref{thm:gap}(2). All layout generators, field offsets
and description sizes can be chosen before installing this
constant-length program. To present $D$ to the evaluator, copy the
active words, including the stored instructions, and fill unused
words with zeros. This costs $O(R)$ ticks. The evaluator generates
pins from these words without any recursive call to the program.

On a successful public path of the simulation, evaluate the
compressor monomial in $\mathbb F_p$ from its public exponent tallies.
Use the same binary-guess test as above to accept with this
multiplicity modulo $p$. Every successful completion of the private
jobs receives the prescribed public weight. The simulation takes
$o(T_-)$ ticks, leaving room for the copy, the prime-field arithmetic
and the extra branch. Thus the program at cap $T_-$ has count
$1+\#E(D)$ modulo $p$. Its genuine wiring has value $Z(D)$, while
the substituted wiring used by $E$ has the very same value.
Figure~\ref{fig:clock} displays the resulting contradiction
\[
                Z(D)=1+\#E(D)=1+Z(D)\quad\text{in }\mathbb F_p .\qedhere
\]
\end{proof}

\begin{figure}[tbp]
\centering
\begin{tikzpicture}[
  font=\small,
  box/.style={draw=ink!65,rounded corners=2pt,align=center,
              inner xsep=6pt,inner ysep=5pt},
  clockstage/.style={box,fill=ink!5,text width=5.2cm,minimum height=1.7cm},
  clockvalue/.style={box,fill=sea!7,draw=sea!70!black,text width=5.2cm},
  arrow/.style={-{Latex[length=2mm]},semithick,draw=ink}]
\node[box,fill=ink!8,text width=9.4cm] (description) at (0,0)
  {One description $D$ stores the batch network and a program\\
   that runs the compressed evaluator on $D$ and adds one.};
\node[clockstage] (run) at (-3.05,-2.2)
  {Run the stored program at $T_-$\\ with its private jobs.};
\node[clockstage] (wire) at (3.05,-2.2)
  {Substitute $R/k$ batch networks\\ in the wiring for $D$. Evaluate\\
   the contraction at $T_+$.};
\node[clockvalue] (runvalue) at (-3.05,-4.1)
  {Value by execution\\[2pt] $Z(D)\equiv 1+\#E(D)$};
\node[clockvalue] (wirevalue) at (3.05,-4.1)
  {Value by substitution\\[2pt] $Z(D)\equiv \#E(D)$};
\draw[arrow] ([xshift=-3.05cm]description.south) --
   node[midway,fill=white,inner sep=2pt]{execute the code} (run.north);
\draw[arrow] ([xshift=3.05cm]description.south) --
   node[midway,fill=white,inner sep=2pt]{read the code as data} (wire.north);
\draw[arrow] (run.south) -- (runvalue.north);
\draw[arrow] (wire.south) -- (wirevalue.north);
\end{tikzpicture}
\caption{Two evaluations of the same pinned contraction in
Theorem~\ref{thm:batches}. The evaluator $E$ on the right counts the
substituted network modulo $p$ for any public code in $D$. Compression
lets the program on the left run that count on its own stored
description and contribute one additional acceptance.}
\label{fig:clock}
\end{figure}

